\documentclass[11pt]{article}
\usepackage{fontspec}
\setmainfont{Times New Roman}
\setsansfont{Arial}
\setmonofont{Consolas}
\usepackage{amsmath,amssymb}
\usepackage{geometry}
\usepackage{microtype}
\geometry{a4paper,margin=3cm}
\setlength{\parindent}{1.25em}
\setlength{\parskip}{0pt}
\emergencystretch=10em
\allowdisplaybreaks
\providecommand{\Kfield}{\mathfrak{K}}
\providecommand{\Ssys}{\mathfrak{S}}
\providecommand{\Jdom}{\mathfrak{J}}
\providecommand{\Gdom}{\mathfrak{G}}
\providecommand{\Lsys}{\mathfrak{L}}
\providecommand{\Omegaint}{[\Omega]}
\providecommand{\eps}{\varepsilon}

\begin{document}

\section*{\S\ 15. Целочислене релативно целе области првого рода и регуларне системы.}

На место интегралној базы при целочисленых системах наступа целочислена интегрална база, то јест конечно число целочисленых полиномов система тако, же всак целочислены полином система се може представити како цела целочислена комбинација сеј конечној множины.

Теоремы о целочисленеј интегралној бази в сучности идут паралелно с предходными, и само доказы требујут мале модификације.

Најпрво дефиниција $\mathrm{V}$ (\S\ 9) може се непосреднье пренести.

{\itshape Дефиниција $\mathrm{V}'$: величина $\xi$ называ се алгебраично цела и целочислена взходно к целочисленеј релативно целеј области $\Gdom_{n\rho}$, ако она задовольа некој једначби, ктореј највышши коефициент јест јединица из $[\Omega]$, а остале полиномы сут из $\Gdom_{n\rho}$.}

Разширенье Гауссова теорема на полиномы с алгебраично-целочислеными коефициентами показује именно, как в \S\ 9, же из того дефиниција може быти добыта на иредуцибилној једначби.

Из дефиниције $\mathrm{V}'$ далье следује пренос дефиниције VII.

{\itshape Дефиниција $\mathrm{VII}'$: целочислена релативно цела област $\Gdom_{n\rho}$ называ се областју првого рода, ако она имаје како област одображеньа релативно целу област $\Gdom^*_{\rho\rho}$ тако, же всак либовольны целочислены полином од $x_1,\ldots,x_\rho$ алгебраично цело и целочислено зависи од $\Gdom^*_{\rho\rho}$.}

Да показати конечност тых областиј првого рода, најпрво заметимо, же как в \S\ 10 из дефиниције следује: највышши коефициент инволуцијској формы јест јединица из $[\Omega]$.

По теорему $\mathrm{VI}'$ (\S\ 14) зато всак полином $F(x)$ из $\Gdom_{n\rho}$ допушча представјенье:
\[
  F(x)=\frac{\Gamma(H_1(x),\ldots,H_\sigma(x))}{i}.
\]
Ако се тут на систем, утворјены из целочисленых полиномов $\Gamma(H_1,\ldots,H_\sigma)$, примени Хилбертов теорем II о целочисленеј модулној бази (Math. Анн. 36), изречены за целе рационалне числове коефициенты, в јего разширеньу на алгебраично-целочислене полиномы\footnote{Ако $w_1,\ldots,w_k$ означа фундаменталны систем алгебраично целых чисел из $\Omega$, и ако поставимо
\[
\begin{gathered}
  \Gamma(H_1,\ldots,H_\sigma)
  =\Gamma_1(H)w_1+\cdots+\Gamma_k(H)w_k,
\end{gathered}
\]
тогда $\Gamma_1(H),\ldots,\Gamma_k(H)$ имајут целе рационалне числове коефициенты. Применьенье теорема II на систем, утворјены из форм $v_1\Gamma_1(H)+\cdots+v_k\Gamma_k(H)$, непосреднье показује важност того разширеньа.}, тогда из того следује обстојанье целочисленеј интегралној базы, то јест:

{\itshape Теорем $\mathrm{VII}'$: всака целочислена релативно цела област првого рода имаје целочислену интегралну базу.}

Далье преносимо дефиницију регуларных системов:

{\itshape Дефиниција $\mathrm{VIII}'$: целочислены систем $\Ssys_{n\rho}$ называ се целочислено-регуларным, ако в области одображеньа $\Jdom^*_{\rho\rho}$ најменшеј содержечеј целочисленеј интегралнеј области обстојат $\rho$ полиномов тако, же резултант членов највышшеј димензије јест јединица из $[\Omega]$, то јест $R_0=\eps$.}

Да пренести доказ обстојаньа интегралној базы, заметимо, же помочно тврđенье из \S\ 11 допушча следујучу острејшу формулацију; она следује из једначб (5) и (6) в \S\ 11 и из околности, же $A_1(\lambda),\ldots,A_\tau(\lambda)$ сут целе целочислене функције од $\lambda_1,\ldots,\lambda_\rho$:

{\itshape Достатно условје за то, же всак либовольны целочислены полином од $x_1,\ldots,x_\rho$ алгебраично цело и целочислено зависи од $\rho$ целочисленых полиномов, јест то, же резултант членов највышшеј димензије тых полиномов јест јединица из $[\Omega]$, то јест $R_0=\eps$.}

Из сеј помочној теоремы следује точно как в \S\ 12 за всак полином $F(x)$ целочислено-регуларного система представјенье (3) из \S\ 12, где ныне $\Gamma_1(G_i),\ldots,\Gamma_k(G_i)$ сут целочислене полиномы од $G_i$. I далье, как в \S\ 12, применьенье Хилбертова теорема II (Math. Анн. 36) в јего разширеньу на алгебраично-целочислене полиномы, вместе с дефиницију најменшеј содержечеј целочисленеј интегралнеј области, даваје обстојанье целочисленеј интегралној базы; то јест:

{\itshape Теорем $\mathrm{IX}'$: всак целочислено-регуларны систем имаје целочислену интегралну базу.}

Как примјер целочисленеј релативно целеј области првого рода, в аналогији с примјером 1 в \S\ 13, указујемо на целост $\Gdom_{nn}$ целочисленых полиномов Лагрангевој родовој области. Же $\Gdom_{nn}$ јест целочислена област првого рода, следује из идентичности:
\[
  x_k^n-\sigma_1(x)x_k^{n-1}+\cdots\pm\sigma_n(x)=0
  \qquad(k=1,2,\ldots,n);
\]
понеже резултант елементарных симетричных функциј имаје ценност $\pm1$, $\Gdom_{nn}$ јест такоже целочислены регуларны систем.

Далејши примјер јест, по помочном теорему из \S\ 14, даны целостју целых целочисленых симетричных функциј од $n$ редов величин.

Ерланген, мај 1914.

\end{document}

